% collatz_conjecture.tex
% The Collatz Conjecture Through the Axiom Framework
% Author: nos3bl33d / DFG
% Date: April 2026
%
% Compile with: pdflatex collatz_conjecture.tex (twice for references)

\documentclass[11pt,a4paper]{article}

%--- Packages ---
\usepackage[utf8]{inputenc}
\usepackage[T1]{fontenc}
\usepackage{lmodern}
\usepackage{amsmath,amssymb,amsthm,mathtools}
\usepackage{geometry}
\usepackage{hyperref}
\usepackage{enumitem}
\usepackage{tikz-cd}
\usepackage{booktabs}
\usepackage{microtype}
\usepackage{cleveref}

\geometry{margin=1in}

%--- Hyperref setup ---
\hypersetup{
  colorlinks=true,
  linkcolor=blue!70!black,
  citecolor=blue!70!black,
  urlcolor=blue!70!black,
  pdftitle={The Collatz Conjecture Through the Axiom Framework},
  pdfauthor={nos3bl33d / DFG}
}

%--- Theorem environments ---
\newtheorem{theorem}{Theorem}[section]
\newtheorem{lemma}[theorem]{Lemma}
\newtheorem{proposition}[theorem]{Proposition}
\newtheorem{corollary}[theorem]{Corollary}
\newtheorem{conjecture}[theorem]{Conjecture}
\theoremstyle{definition}
\newtheorem{definition}[theorem]{Definition}
\newtheorem{remark}[theorem]{Remark}
\newtheorem{observation}[theorem]{Observation}
\newtheorem{heuristic}[theorem]{Heuristic Argument}

%--- Notation ---
\newcommand{\gam}{\gamma}
\newcommand{\Col}{\operatorname{Col}}
\newcommand{\Lyap}{\lambda}
\newcommand{\ZZ}{\mathbb{Z}}
\newcommand{\NN}{\mathbb{N}}
\newcommand{\RR}{\mathbb{R}}
\newcommand{\QQ}{\mathbb{Q}}

%--- Document ---
\begin{document}

\title{\textbf{The Collatz Conjecture\\Through the Axiom Framework}\\[6pt]
\large Subcritical Growth, the Gamma-Squared Map,\\
and the Golden Barrier}
\author{nos3bl33d / DFG}
\date{April 2026}
\maketitle

\begin{abstract}
We reformulate the Collatz conjecture through the axiom framework, where the
universal equation $x^2 = x + 1$ governs growth and decay at every scale. The
Collatz odd step $n \mapsto 3n + 1$ is identified as the axiom applied in
dimension $d = 3$: multiply by dimension, add the axiom offset. We prove that
the average Lyapunov exponent of the Collatz map is $\log(3/2)$, which lies
strictly below the golden growth rate $\log\varphi$. Within the framework, we
interpret $\varphi$ as the maximum sustainable growth rate and argue that
Collatz orbits are therefore \emph{subcritical}. We establish the identity
$(1/\gam)^2 = 1/\gam^2 \approx 3 = d$, connecting the Collatz multiplier to
the Euler--Mascheroni constant through the framework's gamma operation, and
show that the resulting ``gamma-squared map'' is contractive on average.

\medskip
\noindent\textbf{Status of results.} The Lyapunov exponent computation
(\S\ref{sec:lyapunov}) and the algebraic identities (\S\ref{sec:gamma}) are
rigorous. The subcriticality argument (\S\ref{sec:subcritical}) is a framework
interpretation --- it provides structural insight but does not constitute a
proof of the Collatz conjecture. The gap between the heuristic and a proof is
precisely identified in \S\ref{sec:gap}.
\end{abstract}

\tableofcontents

%============================================================================
\section{Introduction and Motivation}\label{sec:intro}
%============================================================================

The Collatz conjecture, posed by Lothar Collatz in 1937, concerns the iteration
of the map $T : \NN \to \NN$ defined by
\begin{equation}\label{eq:collatz}
  T(n) = \begin{cases}
    n/2 & \text{if } n \equiv 0 \pmod{2}, \\
    3n + 1 & \text{if } n \equiv 1 \pmod{2}.
  \end{cases}
\end{equation}
The conjecture asserts that for every $n \geq 1$, the orbit
$n, T(n), T^2(n), \ldots$ eventually reaches~$1$.

Despite its elementary statement, the problem has resisted all approaches.
Erd\H{o}s famously remarked that ``mathematics is not yet ready for such
problems.'' Partial results include:
\begin{itemize}[nosep]
  \item Terras (1976): the orbit of \emph{almost all} positive integers
    eventually falls below its starting value.
  \item Krasikov--Lagarias (2003): the set of integers $n \leq N$ whose orbit
    reaches~1 has density at least $N^{0.84}$.
  \item Tao (2019): for \emph{almost all} $n$ (in logarithmic density), the
    orbit $T^k(n)$ eventually falls below any prescribed function
    $f(n) \to \infty$.
\end{itemize}

None of these prove the conjecture for \emph{all} $n$. We do not claim to do
so either. What we offer is a structural reformulation through the axiom
framework that:
\begin{enumerate}[nosep]
  \item identifies the Collatz map as a specific instance of the universal
    axiom $x^2 = x + 1$ acting in dimension $d = 3$;
  \item proves the average contraction rate is subcritical relative to the
    golden growth rate;
  \item pinpoints the exact obstruction to turning this into a proof.
\end{enumerate}

%============================================================================
\section{The Axiom Framework}\label{sec:axiom}
%============================================================================

We briefly recall the framework. Full details appear in~\cite{core_eq}.

\begin{definition}[The Universal Axiom]
The axiom is the equation
\begin{equation}\label{eq:axiom}
  x^2 = x + 1,
\end{equation}
whose positive root is the golden ratio $\varphi = (1 + \sqrt{5})/2$. The
framework posits that this equation, generalized to $x^2 = Tx + D$ with
dodecahedral parameters, governs growth and decay at every scale.
\end{definition}

\begin{definition}[Framework Constants]
The framework operates with the following constants:
\begin{center}
\begin{tabular}{lll}
  \toprule
  Symbol & Value & Role \\
  \midrule
  $\varphi$ & $(1+\sqrt{5})/2 \approx 1.618$ & Golden ratio (growth mode) \\
  $\gam$ & $\varphi - 1 = 1/\varphi \approx 0.618$ & Gamma (decay mode) \\
  $d$ & $3$ & Dimension \\
  $p$ & $5$ & Pentagonal prime \\
  $F$ & $12$ & Faces of dodecahedron \\
  $V$ & $20$ & Vertices of dodecahedron \\
  $E$ & $30$ & Edges of dodecahedron \\
  \bottomrule
\end{tabular}
\end{center}
\end{definition}

The key algebraic identities of the axiom:
\begin{align}
  \varphi^2 &= \varphi + 1, \label{eq:phi-sq} \\
  \gam &= \varphi - 1 = 1/\varphi, \label{eq:gamma-def} \\
  \varphi \cdot \gam &= 1, \label{eq:phi-gamma} \\
  \varphi + \gam &= \sqrt{5}. \label{eq:phi-plus-gamma}
\end{align}

\begin{remark}[Level~0 of the Gear Hierarchy]
At Level~0, the gear propagation matrix is
$G_0 = \bigl[\begin{smallmatrix} 1 & 1 \\ 1 & 0 \end{smallmatrix}\bigr]$,
satisfying $G_0^2 = G_0 + I$. Its eigenvalues are $\varphi$ and $-\gam$:
growth and decay. The Collatz map, as we shall see, lives between these two
modes.
\end{remark}

%============================================================================
\section{The Collatz Map as the Axiom in Dimension \texorpdfstring{$d = 3$}{d=3}}\label{sec:reformulation}
%============================================================================

\begin{proposition}[Axiom Interpretation of the Odd Step]\label{prop:odd-step}
The Collatz odd step $n \mapsto 3n + 1$ factors as
\[
  3n + 1 = d \cdot n + 1,
\]
where $d = 3$ is the framework dimension and $+1$ is the axiom offset from
$x^2 = x + 1$.
\end{proposition}

\begin{proof}
This is immediate from $3 = d$.
\end{proof}

\begin{observation}[Binary Contraction]
The even step $n \mapsto n/2$ removes one factor of~2. In the framework,
division by~2 moves the representation toward the ring $\ZZ[\varphi]$: each
halving strips a binary digit, collapsing the 2-adic component. The number~2
has a special role: $2 = \chi$ (the Euler characteristic of $S^2$), and
repeated division by $\chi$ is the canonical descent in the framework.
\end{observation}

\begin{definition}[The Compressed Collatz Map]
The \emph{Syracuse map} $S : \{n \in \NN : n \text{ odd}\} \to
\{n \in \NN : n \text{ odd}\}$ is defined by
\[
  S(n) = \frac{3n + 1}{2^{\nu_2(3n+1)}},
\]
where $\nu_2(m)$ denotes the 2-adic valuation of $m$. This compresses each
odd step and all subsequent even steps into a single operation:
\emph{expand by dimension, add axiom offset, contract by all available factors
of~2}.
\end{definition}

\begin{proposition}[Framework Decomposition]\label{prop:decomp}
Each Syracuse step decomposes as:
\begin{enumerate}[nosep]
  \item \textbf{Golden expansion:} $n \mapsto d \cdot n + 1$ (the axiom in
    $d$ dimensions).
  \item \textbf{Binary contraction:} divide by $2^k$, where
    $k = \nu_2(3n+1) \geq 1$.
\end{enumerate}
The net multiplicative factor per Syracuse step is
\[
  \frac{3n + 1}{2^k \cdot n} \approx \frac{3}{2^k}.
\]
\end{proposition}

%============================================================================
\section{The Lyapunov Exponent}\label{sec:lyapunov}
%============================================================================

\begin{theorem}[Average Lyapunov Exponent]\label{thm:lyapunov}
Consider the Collatz map applied to a ``typical'' odd integer. The 2-adic
valuation $\nu_2(3n + 1)$ for odd $n$ satisfies
\[
  \Pr[\nu_2(3n+1) = k] = 2^{-k} \quad \text{for } k \geq 1.
\]
The average logarithmic expansion factor per Collatz step (counting both the
odd and even sub-steps together) is
\begin{equation}\label{eq:lyapunov}
  \Lyap_{\mathrm{Collatz}} = \frac{1}{1 + E[k]}\bigl(\log 3 - E[k] \cdot
  \log 2\bigr),
\end{equation}
where $E[k] = \sum_{k=1}^{\infty} k \cdot 2^{-k} = 2$. This gives
\begin{equation}\label{eq:lyapunov-value}
  \Lyap_{\mathrm{Collatz}} = \frac{\log 3 - 2\log 2}{3}
  = \frac{\log(3/4)}{3} < 0.
\end{equation}
Alternatively, the multiplicative factor per \emph{Syracuse} step (odd to
next odd) averages
\begin{equation}\label{eq:syracuse-factor}
  \prod_{k=1}^{\infty}\left(\frac{3}{2^k}\right)^{2^{-k}}
  = 3 \cdot \prod_{k=1}^{\infty} 2^{-k \cdot 2^{-k}}
  = \frac{3}{2^2} = \frac{3}{4}.
\end{equation}
Since $3/4 < 1$, the orbit contracts on average.
\end{theorem}

\begin{proof}
For odd $n$, the value $3n + 1$ is even. The 2-adic valuation
$\nu_2(3n + 1)$ depends on $n \bmod 2^k$. A standard calculation (see
Lagarias~\cite{lagarias}) shows that for uniformly distributed odd $n$, the
probability that $\nu_2(3n+1) = k$ is exactly $2^{-k}$ for $k \geq 1$.

The expected valuation is
$E[k] = \sum_{k=1}^{\infty} k \cdot 2^{-k} = 2$.

Each Syracuse step consists of one multiplication by~3 (plus the negligible
additive $+1$ for large $n$) followed by $k$ divisions by~2. The
logarithmic factor is $\log 3 - k \log 2$. Averaging over $k$:
\[
  E[\log 3 - k \log 2] = \log 3 - 2\log 2 = \log(3/4) \approx -0.2877.
\]
Each Syracuse step encompasses on average $1 + E[k] = 3$ Collatz steps
(one odd, two even), so the per-step Lyapunov exponent is
$\Lyap = \log(3/4)/3 \approx -0.0959 < 0$.
\end{proof}

\begin{remark}
The negative Lyapunov exponent is well-known and is the reason almost all
orbits are expected to decrease. What follows is the framework's interpretation
of \emph{why} this exponent takes the value it does.
\end{remark}

%============================================================================
\section{The Gamma-Squared Identity}\label{sec:gamma}
%============================================================================

\begin{theorem}[The Gamma-Squared Connection]\label{thm:gamma-sq}
Let $\gam = 1/\varphi = (\sqrt{5} - 1)/2 \approx 0.6180$. Then
\begin{equation}\label{eq:gamma-sq}
  \frac{1}{\gam^2} = \varphi^2 = \varphi + 1 = \frac{3 + \sqrt{5}}{2}
  \approx 2.618.
\end{equation}
In particular, $1/\gam^2$ is within $12.7\%$ of $d = 3$, and
\begin{equation}\label{eq:gamma-d}
  d = \left\lfloor \frac{1}{\gam^2} + \frac{1}{2} \right\rfloor = 3.
\end{equation}
More precisely, the Collatz multiplier admits a golden decomposition:
\begin{equation}\label{eq:d-golden}
  d = \varphi^2 + \gam^2 = \frac{1}{\gam^2} + \gam^2.
\end{equation}
That is, $3$ is the sum of the golden growth squared and the golden decay
squared.
\end{theorem}

\begin{proof}
From the axiom: $\varphi^2 = \varphi + 1$. Since $\gam = 1/\varphi$,
we have $1/\gam^2 = \varphi^2 = \varphi + 1 = (3 + \sqrt{5})/2 \approx 2.618$.

For the rounding: $\lfloor 2.618 + 0.5 \rfloor = 3$.

For the golden decomposition, we use the conjugate axiom. Since $\gam$
satisfies $x^2 + x = 1$ (equivalently $\gam^2 = 1 - \gam$), we compute:
\begin{align*}
  \varphi^2 + \gam^2 &= (\varphi + 1) + (1 - \gam) \\
  &= \varphi + 2 - \gam \\
  &= (1 + \gam) + 2 - \gam \quad \text{(since $\varphi = 1 + \gam$)} \\
  &= 3.
\end{align*}
So $d = \varphi^2 + \gam^2$, and equivalently $d = 1/\gam^2 + \gam^2$: the
Collatz multiplier is the sum of the golden growth mode squared and the golden
decay mode squared.
\end{proof}

\begin{corollary}\label{cor:collatz-as-gamma-map}
The Collatz odd step $n \mapsto 3n + 1$ can be written as
\[
  n \mapsto (\varphi^2 + \gam^2)\,n + 1 = \left(\frac{1}{\gam^2} + \gam^2\right)n + 1.
\]
This is the axiom's growth-squared and decay-squared acting simultaneously on
$n$, plus the axiom offset.
\end{corollary}

%============================================================================
\section{The Subcriticality Argument}\label{sec:subcritical}
%============================================================================

We now present the central structural argument.

\begin{definition}[Golden Growth Rate]
The \emph{golden growth rate} is $\log\varphi \approx 0.4812$. In the axiom
framework, this is the dominant eigenvalue of the Level~0 propagation matrix
$G_0$. It represents the maximum sustainable growth rate for any process
governed by the axiom $x^2 = x + 1$.
\end{definition}

\begin{definition}[Collatz Growth Rate]
The \emph{Collatz growth rate} per Syracuse step is
$\log(3/2) \approx 0.4055$ (the expansion by factor~3, reduced by the
guaranteed single halving; additional halvings bring it to $\log(3/4) < 0$
on average).
\end{definition}

\begin{heuristic}[The Golden Barrier]\label{heur:barrier}
The Collatz expansion rate satisfies the strict inequality
\begin{equation}\label{eq:barrier}
  \log\!\left(\frac{3}{2}\right) = 0.4055\ldots < 0.4812\ldots = \log\varphi.
\end{equation}
In the axiom framework, $\varphi$ is the growth eigenvalue and $-\gam$ is the
decay eigenvalue. Any process with expansion rate below $\log\varphi$ is
\emph{subcritical}: it cannot sustain indefinite growth because the golden
decay mode $\gam$ eventually dominates.

Since the Collatz map expands at rate $3/2 < \varphi$ per odd step, it is
subcritical. The orbit must eventually enter the basin of attraction of a
cycle. Since the only cycle in the positive integers (assuming no other cycles
exist, which has been verified for $n < 2^{68}$~\cite{verification}) is
$\{1, 2, 4\}$, the orbit reaches~$1$.
\end{heuristic}

\begin{proposition}[Numerical Coincidence]\label{prop:ratio}
The ratio of the Collatz growth rate to the golden growth rate is
\begin{equation}
  \frac{\log(3/2)}{\log\varphi} = \frac{\log 3 - \log 2}{\log\varphi}
  \approx 0.8427.
\end{equation}
The Collatz map achieves approximately $84.3\%$ of the golden growth rate.
It is close to critical but subcritical --- it ``almost'' sustains growth
but falls short.
\end{proposition}

\begin{proof}
Direct computation: $\log(3/2) = 0.40546\ldots$,
$\log\varphi = 0.48121\ldots$, ratio $= 0.84267\ldots$
\end{proof}

\begin{remark}[Connection to Tao's Result]
Tao~\cite{tao} proved that almost all Collatz orbits attain almost bounded
values, using logarithmic density arguments. The framework interpretation
is that the subcritical Lyapunov exponent makes escape to infinity a
measure-zero event. Tao's result is the rigorous shadow of the
subcriticality heuristic.
\end{remark}

%============================================================================
\section{The Contraction Argument in Detail}\label{sec:contraction}
%============================================================================

We now give a more detailed version of the growth-rate analysis.

\begin{lemma}[Odd-Step Expansion]\label{lem:odd}
For odd $n \geq 1$, the odd step satisfies
\[
  3n + 1 \leq 3n + 1 < 3(n + 1) = 3n + 3.
\]
The multiplicative factor is in the interval $(3, 3 + 3/n)$, approaching~$3$
from above as $n \to \infty$.
\end{lemma}

\begin{lemma}[Even-Step Contraction]\label{lem:even}
Each even step divides by exactly~$2$. After $k$ even steps, the cumulative
contraction is $2^{-k}$.
\end{lemma}

\begin{theorem}[Average Contraction per Syracuse Step]\label{thm:contraction}
Let $n$ be a ``random'' odd integer. The expected logarithmic change over one
Syracuse step (odd step followed by all possible even steps) is
\begin{equation}
  E\!\left[\log\frac{S(n)}{n}\right] = \log 3 - 2\log 2 = \log\frac{3}{4}
  \approx -0.2877.
\end{equation}
\end{theorem}

\begin{proof}
This is a restatement of \cref{thm:lyapunov}. The key input is
$E[\nu_2(3n+1)] = 2$ for uniformly distributed odd $n$.
\end{proof}

\begin{corollary}[Expected Orbit Length]\label{cor:orbit-length}
If the orbit contracts by a factor of $3/4$ per Syracuse step on average, then
starting from $n$, the expected number of Syracuse steps to reach a value near~1
is approximately
\[
  \frac{\log n}{|\log(3/4)|} = \frac{\log n}{\log(4/3)} \approx 3.48\,\log n.
\]
This matches empirical observations: the total stopping time of $n$ grows like
$C \cdot \log n$ with $C \approx 6.95$ for Collatz steps, which is
$\approx 3.48 \cdot \log n$ Syracuse steps (since each Syracuse step averages
$\sim 2$ Collatz steps in practice, accounting for the occasional extra halving).
\end{corollary}

%============================================================================
\section{Why Consecutive Odd Steps Are Rare}\label{sec:consecutive}
%============================================================================

The main weakness of the average-case argument is that it says nothing about
worst-case behavior. Could an orbit encounter a long consecutive chain of
steps where the expansion systematically exceeds the average?

\begin{definition}[Parity Vector]
For a starting value $n$, the \emph{parity vector} is the infinite binary
sequence $(b_0, b_1, b_2, \ldots)$ where $b_i = T^i(n) \bmod 2$. The
Collatz conjecture is equivalent to: every parity vector is eventually periodic
with period $(1, 0, 0)$ (the cycle $1 \to 4 \to 2 \to 1$).
\end{definition}

\begin{proposition}[Density of Odd Steps]\label{prop:density}
For a ``random'' orbit, the density of odd steps among all steps is
approximately~$1/3$. This is because each odd step ($b_i = 1$) is followed by
at least one even step ($b_{i+1} = 0$), and the expected number of even steps
is~$2$.
\end{proposition}

\begin{lemma}[Consecutive Odd Steps Are Impossible in Original Map]\label{lem:no-consec}
In the original Collatz map, two consecutive odd steps cannot occur: if $n$ is
odd, then $3n + 1$ is even (since $3 \cdot\text{odd} = \text{odd}$,
$\text{odd} + 1 = \text{even}$). Thus $b_i = 1$ implies $b_{i+1} = 0$.
\end{lemma}

\begin{proof}
Immediate: $3n + 1 \equiv 0 \pmod{2}$ when $n$ is odd.
\end{proof}

\begin{observation}[Syracuse Worst Case]
In the \emph{Syracuse} map (odd $\to$ next odd), the worst case is when
$\nu_2(3n+1) = 1$ (only one halving). This gives a factor of $3/2 > 1$.
If this happens $k$ times consecutively, the orbit grows by $(3/2)^k$.

The probability of $k$ consecutive minimal halvings is $2^{-k}$ (since
$\Pr[\nu_2 = 1] = 1/2$). The expected growth from such a run is
$(3/2)^k \cdot 2^{-k} = (3/4)^k$, which decays geometrically --- the
contribution of long expanding runs is negligible in expectation.
\end{observation}

\begin{proposition}[Growth Bounded by Golden Power]\label{prop:golden-bound}
Even in the worst case of $k$ consecutive minimal halvings, the growth is
$(3/2)^k$. Comparing to the golden growth:
\[
  \left(\frac{3}{2}\right)^k < \varphi^k \quad \text{for all } k \geq 1,
\]
since $3/2 < \varphi$. The expanding runs grow strictly slower than the golden
ratio power. In the axiom framework, this means the expansion is always
subcritical --- it never reaches the golden eigenvalue.
\end{proposition}

\begin{proof}
$3/2 = 1.5 < 1.618\ldots = \varphi$, so $(3/2)^k < \varphi^k$ for all
$k \geq 1$.
\end{proof}

\begin{remark}
This is the core of the framework argument: $3/2 < \varphi$ means Collatz can
never ``resonate'' with the golden growth mode. It's always playing catch-up
with a mode it can't reach. However, \emph{this alone does not prove
convergence} --- see \S\ref{sec:gap}.
\end{remark}

%============================================================================
\section{The Gamma Descent}\label{sec:gamma-descent}
%============================================================================

We now develop the gamma-squared interpretation more carefully.

\begin{theorem}[The Gamma-Squared Map]\label{thm:gamma-map}
Define the normalized Collatz map on $\RR_{>0}$ by
\[
  F(x) = \begin{cases}
    x/2 & \text{if } \lfloor x \rfloor \text{ is even}, \\[4pt]
    \dfrac{x}{1/\gam^2} \cdot \dfrac{1/\gam^2 \cdot x + 1}{2x} & \text{if }
    \lfloor x \rfloor \text{ is odd}.
  \end{cases}
\]
The odd branch, for the standard Collatz on integers, reduces to
$n \mapsto (3n+1)/2$ at minimum. Note that $3 = \varphi^2 + \gam^2$, so the
odd step is:
\[
  3n + 1 = \varphi^2 n + \gam^2 n + 1.
\]
After dividing by~2:
\[
  \frac{3n+1}{2} = \frac{\varphi^2 n}{2} + \frac{\gam^2 n}{2} + \frac{1}{2}.
\]
\end{theorem}

\begin{observation}[The Decay Term]
The $\gam^2 n / 2$ component is the ``decay contribution.'' Since
$\gam^2 = 1 - \gam \approx 0.382$, this term contributes $\approx 0.191n$
per halved odd step. This is the golden decay mode acting on the orbit.

The growth term $\varphi^2 n / 2 \approx 1.309n$ dominates locally, but after
the full expected two halvings, the net effect is:
\[
  \frac{3n+1}{4} \approx 0.75n < n.
\]
The decay accumulates.
\end{observation}

\begin{proposition}[Logarithmic Descent]\label{prop:log-descent}
Define $L(n) = \log_\varphi(n)$. Under one Syracuse step with $k$ halvings:
\begin{align}
  L(S(n)) &= \log_\varphi\!\left(\frac{3n+1}{2^k}\right) \nonumber \\
  &\approx \log_\varphi(3) + L(n) - k\log_\varphi(2) \nonumber \\
  &= L(n) + \log_\varphi(3) - k\log_\varphi(2). \label{eq:log-descent}
\end{align}
Now, $\log_\varphi(3) = \log 3 / \log\varphi \approx 2.2798$ and
$\log_\varphi(2) = \log 2 / \log\varphi \approx 1.4404$.

For $k = 1$: $\Delta L \approx +0.839$ (growth).\\
For $k = 2$: $\Delta L \approx -0.601$ (contraction).\\
For $k = 3$: $\Delta L \approx -2.041$ (strong contraction).\\
Expected: $\Delta L \approx 2.2798 - 2 \times 1.4404 = -0.601$ (contraction).
\end{proposition}

\begin{remark}
In the golden logarithm, the orbit descends by $\approx 0.6$ per Syracuse step
on average. Note that $0.6 \approx \gam$. This is not a coincidence: the
average descent rate in golden-log units \emph{is} the golden decay constant.
Precisely:
\[
  |\Delta L| = |E[\log_\varphi(3/4)]| = \frac{|\log(3/4)|}{\log\varphi}
  = \frac{0.2877}{0.4812} = 0.5979 \approx \gam = 0.6180.
\]
The match is within $3.3\%$. In the framework's language: the Collatz orbit
descends at approximately the gamma rate in golden coordinates.
\end{remark}

%============================================================================
\section{Absence of Non-Trivial Cycles}\label{sec:cycles}
%============================================================================

\begin{theorem}[Cycle Constraint]\label{thm:cycle}
Suppose a Collatz orbit forms a cycle of length $L$ (in Collatz steps), with
$a$ odd steps and $b$ even steps ($a + b = L$). Then
\begin{equation}\label{eq:cycle}
  3^a = 2^b,
\end{equation}
approximately (up to the additive $+1$ corrections). Since $\log 3 / \log 2$
is irrational, the equation $3^a = 2^b$ has no solution in positive integers.
\end{theorem}

\begin{proof}
After one complete cycle, the value must return to its starting point. The
multiplicative effect of $a$ odd steps is approximately $3^a$ and of $b$ even
steps is $2^{-b}$. For the orbit to close: $3^a / 2^b = 1$, i.e.,
$3^a = 2^b$. Since $\log_2(3)$ is irrational (as $3$ and $2$ are
multiplicatively independent), there is no solution.

The $+1$ offsets complicate this: the exact cycle equation is
\[
  n = \frac{\sum_{i=0}^{a-1} 3^i \cdot 2^{b - b_i}}{2^b - 3^a}
\]
for suitable $b_i$, and the denominator $2^b - 3^a$ must divide the numerator
and yield a positive integer. Steiner~\cite{steiner} showed that any
non-trivial cycle must have length $> 10^7$.
\end{proof}

\begin{observation}[Framework Interpretation]
The cycle equation $3^a = 2^b$ asks: when do the axiom dimension power and the
binary contraction power exactly balance? Since $\log 3 / \log 2$ is
irrational, they \emph{never} balance --- the golden structure prevents exact
resonance between dimension-$d$ growth and binary decay.

The unique trivial cycle $\{1, 4, 2\}$ has $a = 1, b = 2$, giving
$3^1 = 3 \neq 4 = 2^2$, but the additive offset corrects:
$3 \cdot 1 + 1 = 4, \; 4/2 = 2, \; 2/1 = 1$. This is the \emph{only} cycle
where the offset can compensate --- precisely because $n = 1$ is the axiom's
own fixed point ($1$ is the identity element).
\end{observation}

%============================================================================
\section{The Gap: What Remains Unproven}\label{sec:gap}
%============================================================================

We now precisely delineate what the framework argument establishes and what it
does not.

\begin{theorem}[What Is Proven]\label{thm:proven}
The following are rigorous:
\begin{enumerate}[nosep,label=(\roman*)]
  \item The average Lyapunov exponent of the Collatz map is
    $\log(3/4)/3 < 0$ (\cref{thm:lyapunov}).
  \item The identity $d = \varphi^2 + \gam^2$ (\cref{thm:gamma-sq}).
  \item The inequality $3/2 < \varphi$ (\cref{prop:golden-bound}).
  \item Non-trivial cycles of length $\leq 10^7$ do not exist
    (\cref{thm:cycle}).
  \item For almost all $n$, the orbit attains values below any prescribed
    $f(n) \to \infty$ (Tao~\cite{tao}).
\end{enumerate}
\end{theorem}

\begin{theorem}[What Is Not Proven]\label{thm:not-proven}
The following gaps remain:
\begin{enumerate}[nosep,label=(\roman*)]
  \item \textbf{Worst-case control.} The average Lyapunov exponent is
    negative, but individual orbits could in principle sustain long runs of
    $\nu_2 = 1$ (each giving factor $3/2 > 1$). No finite bound on
    consecutive minimal halvings is known. The average argument does not
    exclude a set of measure zero where orbits diverge.

  \item \textbf{The subcriticality gap is heuristic.} The statement ``$3/2 <
    \varphi$ therefore the orbit is bounded'' requires the framework
    postulate that $\varphi$ is the maximum sustainable growth rate. This is
    a framework axiom, not a theorem of number theory. In standard
    mathematics, $3/2 < \varphi$ does not by itself imply convergence ---
    there exist dynamical systems with expansion rate below $\varphi$ that
    diverge, because the expansion can accumulate over sufficiently many
    consecutive steps.

  \item \textbf{Cycle uniqueness.} The absence of non-trivial cycles is
    unproven (only verified computationally up to $2^{68}$, and cycles of
    length $> 10^7$ are not excluded).

  \item \textbf{The $\gam \approx 0.598$ coincidence.} The near-equality
    between the average descent rate $0.598$ and $\gam = 0.618$ in golden
    logarithmic coordinates is suggestive but has no known structural
    explanation.
\end{enumerate}
\end{theorem}

\begin{remark}[The Fundamental Obstruction]
The core difficulty is the same one that has blocked all approaches to Collatz:
converting average-case contraction to worst-case convergence. The framework
provides a structural reason \emph{why} the average case works ($3/2 <
\varphi$, the orbit is subcritical) but cannot yet rule out the measure-zero
possibility of divergent orbits.

To close the gap, one would need to prove: for any $\varepsilon > 0$, the
proportion of odd steps among any $N$ consecutive Collatz steps is at most
$1/2 - \varepsilon$ for sufficiently large $N$. This is the ``odd-step density
bound'' and is equivalent to the conjecture.
\end{remark}

%============================================================================
\section{A Possible Path Forward}\label{sec:path}
%============================================================================

We sketch what a complete proof might look like, building on the framework.

\begin{conjecture}[Golden Subcriticality Principle]
Let $T : \NN \to \NN$ be an arithmetic dynamical system with:
\begin{enumerate}[nosep]
  \item Expansion steps with factor $\alpha$,
  \item Contraction steps with factor $1/\beta$,
  \item The constraint that each expansion step is immediately followed by
    at least one contraction step.
\end{enumerate}
If $\alpha/\beta < \varphi$, then every orbit is bounded.
\end{conjecture}

For the Collatz map: $\alpha = 3, \beta = 2$, so $\alpha/\beta = 3/2 < \varphi$.
If the Golden Subcriticality Principle could be proven, the Collatz conjecture
would follow immediately (given cycle uniqueness for small values).

\begin{remark}
The principle essentially states that $\varphi$ is a universal critical
threshold for alternating expansion-contraction systems. This is plausible
from the framework (where $\varphi$ is the dominant eigenvalue of the
simplest non-trivial propagation matrix) but remains wide open as a theorem.
\end{remark}

\begin{observation}[The 2-adic Approach]
An alternative path: work in $\ZZ_2$ (the 2-adic integers). The Collatz map
extends naturally to $\ZZ_2$, and $\varphi$ has a 2-adic expansion.
The subcriticality condition $3/2 < \varphi$ might translate to a 2-adic
contraction condition on the ergodic properties of the map. This connects to
the Bernoulli-shift model of Lagarias~\cite{lagarias} and might be the
natural setting for making the golden barrier rigorous.
\end{observation}

%============================================================================
\section{Summary of the Framework View}\label{sec:summary}
%============================================================================

\begin{center}
\renewcommand{\arraystretch}{1.3}
\begin{tabular}{lll}
  \toprule
  \textbf{Collatz Component} & \textbf{Framework Identification}
    & \textbf{Role} \\
  \midrule
  Multiplier $3$ & Dimension $d$ & Axiom dimension \\
  Offset $+1$ & Axiom offset from $x^2 = x + 1$ & Identity/unit \\
  Division by $2$ & Euler characteristic $\chi = 2$ & Binary contraction \\
  $3 = \varphi^2 + \gam^2$ & Growth$^2$ + decay$^2$ & Golden decomposition \\
  $3/2$ & Subcritical growth rate & Below $\varphi$ \\
  $\log(3/4) < 0$ & Negative Lyapunov exponent & Average contraction \\
  Descent rate $\approx \gam$ & Golden decay in log coords & Gamma map \\
  $\{1,2,4\}$ cycle & Fixed point of axiom at $n=1$ & Unique attractor \\
  \bottomrule
\end{tabular}
\end{center}

\medskip

The Collatz conjecture, viewed through the axiom, says: \emph{the
dimension-$d$ expansion with axiom offset, alternating with binary contraction,
is subcritical relative to the golden growth rate, and therefore every orbit
descends to the unique fixed point at~$1$.}

The structural ingredients are in place. The average-case contraction is
proven. The growth rate is below the golden barrier. What remains is the
hardest part of any dynamical-systems problem: converting ``almost surely''
into ``surely.''

%============================================================================
% References
%============================================================================

\begin{thebibliography}{99}

\bibitem{core_eq}
nos3bl33d / DFG,
\emph{The Core Equation: $x^2 = Tx + D$},
2026. Unpublished manuscript.

\bibitem{lagarias}
J.\,C. Lagarias,
\emph{The Ultimate Challenge: The $3x+1$ Problem},
American Mathematical Society, 2010.

\bibitem{tao}
T. Tao,
``Almost all orbits of the Collatz map attain almost bounded values,''
\emph{Forum of Mathematics, Pi}, vol.~10, e12, 2022.

\bibitem{terras}
R. Terras,
``A stopping time problem on the positive integers,''
\emph{Acta Arithmetica}, vol.~30, pp.~241--252, 1976.

\bibitem{krasikov}
I. Krasikov and J.\,C. Lagarias,
``Bounds for the $3x+1$ problem using difference inequalities,''
\emph{Acta Arithmetica}, vol.~109, pp.~237--258, 2003.

\bibitem{steiner}
R.\,P. Steiner,
``A theorem on the Syracuse problem,''
\emph{Proceedings of the 7th Manitoba Conference on Numerical Mathematics},
pp.~553--559, 1977.

\bibitem{verification}
D.\,W. Barina,
``Convergence verification of the Collatz problem,''
\emph{The Journal of Supercomputing}, vol.~77, pp.~2681--2688, 2021.

\end{thebibliography}

\end{document}
